\section{The retained stationary region and its exterior}
\label{sec:stationary-exterior-receipt}
The Gaussian comparison must describe a region of the actual fine
configuration space. In particular, the fine-field cutoff must be
one throughout that region. We establish this using the funded
expansion to find a configuration with slack, then propagating the
slack by a uniform Lipschitz estimate.

\subsection{Coordinates and the normalized integral}
Fix an auxiliary finite torus, a tree gauge, and the actual coarse
field $V(x)$ in logarithmic coordinates $x\in X$. Let $Y$ be the
Euclidean space of free fine logarithmic coordinates, $d=\dim Y$,
and $N$ the number of free fine bonds. Its coordinate map is
\[
 {\cal D}_x(y)_i=\exp(y_i)\,V(x)^{\rm rep}_i.
\]
Here $y_i$ includes the fixed orthonormal identification with the
Lie algebra, and tree bonds are restored by the gauge fixing. The space $X$
uses the supremum norm on real color coordinates. The space $Y$
uses its Euclidean norm for Gaussian integration; write
$\|y\|_\infty$ for the supremum norm after forgetting that
Euclidean norm. Thus $\|y\|_\infty\le\|y\|_2$. The combined
chart norm is $\max(\|x\|_\infty,\|y\|_\infty)$.
Write $\rho(x)$ for the actual stationary branch, $b=\beta^{-1/2}$,
and
\[
 \mathcal S_x={\cal D}_x\bigl(\rho(x)+b\overline B_Y(0,D)\bigr).
\]
This is a compact measurable subset of the free configuration
space. All Haar measures in this section are probability measures.

The local Haar density of Section~\ref{sec:haar-chart} is
$c_H J(y_i)$ on each free bond. Set
\[
 \psi(x,y)=\ell({\cal D}_x(y))-\sum_{i=1}^{N}\log J(y_i),
 \qquad M(x)=D_y^2\phi(x,\rho(x)),
\]
where $\phi$ is the fine classical action plus the smoothing
penalty, and $\ell$ is the fine loop correction. Assume that the
stationary chart estimates hold on the region in question, with
$M(x)>0$. Introduce
\[
 \Phi(x)=\beta\phi(x,\rho(x))+\psi(x,\rho(x))
                         +\tfrac12\log\det M(x),\qquad
 \Omega=\pi_\beta(\sqrt{2\pi}\,b)^d c_H^N,
\]
where $\pi_\beta>0$ is the smoothing-kernel prefactor.
The tagged integrand of the pure push is
\begin{equation}\label{eq:exterior-tag-integrand}
 U_x(U)=\sum_k\pi_\beta\lambda_k(U)
   e^{-\beta[ a(U)+n_k(U,V(x))]-\ell(U)}
                    {\bf1}_{\{\text{fine field healthy}\}}.
\end{equation}
The nonnegative tag weights satisfy $\sum_k\lambda_k(U)=1$;
a dead tag has zero weight. Tree-gauge integration identifies the
pure push with $Z(x)=\int U_x\,dU$. The incoming density and kernel
have integrable envelopes, so all restricted integrals below exist.

Suppose every point of $\mathcal S_x$ is fine healthy, the whole
logarithmic ball is inside the Haar chart, and all live tag actions
there equal $\phi$. Then the product chart formula and the affine
change $y=\rho(x)+bz$ give
\begin{equation}\label{eq:exterior-exact-interior}
       \int_{\mathcal S_x}U_x\,dU
                  =\Omega e^{-\Phi(x)}R(b,x),
\end{equation}
with exactly the truncated integral $R$ of
Section~\ref{sec:truncated-stationary-log}.
Indeed the Jacobian is $b^dc_H^N\prod_iJ(y_i)$, the tag sum is
one, and $\beta b^2=1$. The quadratic Gaussian factor contributes
$(2\pi)^{d/2}(\det M(0))^{-1/2}$; its determinant and the relative
determinant in the exponent defining $R$ combine to
$(\det M(x))^{-1/2}$. The remaining exponent is precisely
$-\Phi(x)$. No tag-count factor occurs. Injectivity of the Haar
chart ensures that the inverse image of $\mathcal S_x$ within that
chart is the chosen closed ball, so this change of variables also
handles its boundary.

\subsection{The energy cost outside the retained region}
Let $\gamma>0$ be the stationary chart gap and put
$\Delta=\gamma D^2/4$. On the common safe chart assume the fine
Hessian satisfies $D_y^2\phi\ge(\gamma/2)I$.
Stationarity and Taylor's integral formula then imply
\begin{equation}\label{eq:exterior-convex-gap}
 \phi(x,y)-\phi(x,\rho(x))
          \ge\frac\gamma4\|y-\rho(x)\|^2.
\end{equation}
Here the safe chart is the inverse image of the closed combined
coordinate supremum ball, while the displayed quadratic lower
bound uses the Euclidean norm on $Y$. Both endpoints lie in that
chart, and its convexity contains their entire joining segment.
Consequently any point of that chart outside $\mathcal S_x$ pays
at least $\Delta$ after multiplication by $\beta$.

We also need to recognize chart points without presupposing that
the fine field is in the chart. The tree-gauge geometric estimate
supplies a constant $C_0$ such that an alive tag, fine plaquette
distances at most $\delta$, and prediction-wire distances at most
$\delta$ imply
\[
       \operatorname{dist}(U_i,V(x)^{\rm rep}_i)
                       \le (C_0s+1)\delta
\]
for every free bond. Here $s$ bounds the number of plaquettes in
the fixed local comparison box. The estimate follows by comparing
that bond with its tree-gauge replica: the local filling contributes
at most $s\delta$, its geometric comparison costs $C_0$, and the
prediction wire contributes the final $\delta$.
Choose logarithm-chart constants $r_0,C_{\log}>0$, so that distances
at most $r_0$ have logarithms bounded by $C_{\log}$ times the
distance. For a prescribed safe coordinate radius $r_s>0$, choose
\[
 \delta\le\min\left\{\frac{r_0}{C_0s+1},
                    \frac{r_s}{C_{\log}(C_0s+1)}\right\}.
\]
Taking the logarithm of each relative free bond now yields
$U={\cal D}_x(y)$ with $\|y\|_\infty\le r_s$. The coarse
coordinates are already inside the safe radius. This argument uses
the geometric comparison for alive tags only.

Write $\tau$ for the fine health threshold, and suppose
$\tau\le\delta$. Let $W\ge1$ bound all relevant site and coordinate
counts. The path budgets on a healthy field give
\[
 |\ell(U)|\le W\Lambda\tau^2,\qquad
 a(U)\ge-WC(\tau^3+\tau^4).
\]
The second bound uses the nonnegative quadratic part of the action
and the absolute cubic and quartic budgets. Define
\[
 B_{\rm loop}=\beta C(\tau^3+\tau^4)+\Lambda\tau^2.
\]
Let $m_W$ bound $|\phi(x,\rho(x))|/\|x\|$ on the coarse chart,
and let $p+t_W$ bound
$|\Phi(x)-\beta\phi(x,\rho(x))|$. The latter consists of the
loop/Haar bound $p$ and a log-determinant bound $t_W$.
If a prediction wire exceeds $\delta$, its nonnegative smoothing
penalty is at least $\kappa\delta^2/2$, where $\kappa>0$ is the
kernel penalty constant. Assume
\begin{equation}\label{eq:exterior-wire-reserve}
            \beta m_W\|x\|+\Delta
                              \le\beta\kappa\delta^2/2.
\end{equation}
For a healthy point outside $\mathcal S_x$, there are now two cases.
If all prediction wires are small, the preceding chart-recognition
argument and \eqref{eq:exterior-convex-gap} apply. Otherwise the
large-wire penalty and \eqref{eq:exterior-wire-reserve} apply.
In both cases the exponent in each live tag satisfies
\[
 -\beta[a(U)+n_k(U,V(x))]-\ell(U)
           \le-\Phi(x)-\Delta+WB_{\rm loop}+p+t_W.
\]
Unhealthy fields and dead tags contribute zero. Summing the tag
weights therefore proves a pointwise exterior bound.

To express it in stationary units, put
$m_H=c_H(\sqrt{2\pi}\,b)^{\dim\mathfrak g}$, so
$\Omega=\pi_\beta m_H^N$, and define
\begin{equation}\label{eq:exterior-exponent}
       Y_*=-N\log m_H-\Delta+WB_{\rm loop}+p+t_W.
\end{equation}
Then
\begin{equation}\label{eq:exterior-pointwise}
 U_x(U)\le\Omega e^{-\Phi(x)}e^{Y_*}
                              \quad(U\notin\mathcal S_x).
\end{equation}
Since the ambient Haar measure has mass one, the same bound holds
for the exterior integral.

\subsection{Finding one configuration with slack}
Let $Z_{\rm sl}(x)$ denote the portion of the pure push allocated
to the empty coarse and fine defect sector. The sector allocation
has values between zero and one. Its nonzero support requires
all fine plaquette distances to be at most $\tau/2$.
The exact funded expansion of this sector is
\[
        Z_{\rm sl}(x)={\cal O}(x)\sum_I e_I(x),
        \qquad e_\varnothing(x)=1,
\]
where ${\cal O}(x)>0$ is the same sector overhead as in the block
integral, and the expansion gives
$\sum_I|e_I(x)|\le e^\varepsilon$. These are the actual funded
sector identities, with its fixed tree and allocation. They imply
\begin{equation}\label{eq:exterior-slack-lower}
 Z_{\rm sl}(x)\ge {\cal O}(x)(2-e^\varepsilon).
\end{equation}
To see the factor two, separate the empty term: the absolute sum
of the remaining terms is at most $e^\varepsilon-1$, so their sum
is at least $1-e^\varepsilon$. The expansion's absolute convergence
justifies this separation. On the auxiliary volume, its rooted
budget gives $\varepsilon=|\Lambda_c|\,o(x)$, with a negative upper
exponent for $o$. Since $|\Lambda_c|$ has upper exponent zero,
$e^\varepsilon\le3/2$ eventually.

We compare ${\cal O}$ to the stationary normalization before using
\eqref{eq:exterior-slack-lower}. On a healthy coarse field its exact
formula is
\[
       {\cal O}=\Omega
          \exp\{-\beta a_{{\rm coarse},2}(x)
                            -\tfrac12\log\det A_0\}.
\]
Here $a_{{\rm coarse},2}$ is the coarse quadratic action and
$A_0$ is the empty-sector covariance precision matrix. The funded
matrix bound gives $\log\det A_0\le W a_*$. Classical matching
through degree four, with fifth-derivative cost $K_5$, and the
coarse cubic and quartic budgets give
\[
 |\phi(x,\rho(x))-a_{\rm coarse}(x)|\le WK_5\|x\|^5,
 \quad |a_{\rm coarse}(x)-a_{{\rm coarse},2}(x)|
                         \le WC'(\tau_c^3+\tau_c^4).
\]
Combining these inequalities with the bound $p+t_W$ yields
\begin{equation}\label{eq:exterior-overhead-compare}
 {\cal O}\ge\Omega e^{-\Phi-e_*},\qquad
 e_*=\beta W[K_5\|x\|^5+C'(\tau_c^3+\tau_c^4)]
                                    +p+t_W+Wa_*/2.
\end{equation}
If
\begin{equation}\label{eq:exterior-two-reserves}
       Y_*+e_*+2\le-\Delta/2,\qquad Y_*\le-\Delta/2,
\end{equation}
then the exterior upper bound is strictly less than
$\tfrac12{\cal O}\le Z_{\rm sl}$, since $e^{-2}<1/2$.
If $\mathcal S_x$ contained no point with all fine plaquettes at
most $\tau/2$, the allocated integrand would vanish throughout
$\mathcal S_x$. Outside it, that integrand is bounded by $U_x$.
Its integral would consequently be bounded by the exterior upper
bound, a contradiction. Thus there exists
\[
 U_0={\cal D}_x(\rho(x)+bz_0),\qquad \|z_0\|\le D,
        \qquad\operatorname{toll}_i(U_0)\le\tau/2\quad\hbox{for all }i.
\]
This existence argument does not differentiate the fine cutoff.

\subsection{Spreading the slack and choosing the parameters}
The fine plaquette tolls are uniformly Lipschitz, with a constant
$L_0>0$ independent of the auxiliary volume, on the combined
coordinate supremum unit ball, with Lipschitz distance measured
in that same norm. Indeed each plaquette contains four links;
the uniform first-derivative estimate for a link, the product rule,
and unitarity bound the derivative of its matrix holonomy. The
reverse triangle inequality for the matrix norm then bounds the
difference of tolls. Integrating that derivative along the joining
segment proves the asserted Lipschitz estimate.

Let $C_\rho$ be the supremum-coordinate Lipschitz bound for the stationary
branch, with $\rho(0)=0$. For $\|y-x\|\le r$ and $\|z\|\le D$,
the supremum distance between the combined coordinates of
${\cal D}_y(\rho(y)+bz)$ and $U_0$ is at most
\[
                       (1+C_\rho)r+2bD.
\]
We used $\|z-z_0\|_\infty\le\|z-z_0\|_2\le2D$.
Thus every such field is fine healthy provided
\begin{equation}\label{eq:exterior-slack-propagation}
                  L_0[(1+C_\rho)r+2bD]\le\tau/2.
\end{equation}
The two configurations lie in the coordinate unit ball under the
radius conditions below. The same conditions place the entire
region inside the Haar chart and inside the region where all tag
actions coincide.

Here are the simultaneous purchases used in the construction.
The action constants $C,C'$ are fixed and nonnegative, while
$\kappa,\gamma$ are fixed and positive. The constant $C_\rho$
is fixed at this stage independently of the auxiliary volume.
The loop budget obeys $0\le\Lambda\le2P_F$.
Write $z=\beta_F$, $P_F$ for the fine monitor,
$u=z^{-7/16}$ and $r=z^{-1}$, and take
\[
 \beta\in[z/2,2z],\qquad \tau=z^{-1/2}P_F,\qquad
       D=d_0P_F,\qquad d_0=(16L_0)^{-1}.
\]
Every fixed polynomial in the auxiliary size parameter is
$o(P_F^\epsilon)$ for every $\epsilon>0$; $P_F$ itself has upper
exponent zero. These are the monitor's two distinct growth
properties. Since $b\le2z^{-1/2}$,
$2L_0bD\le\tau/4$, and
$L_0(1+C_\rho)r\le\tau/4$ eventually. This proves
\eqref{eq:exterior-slack-propagation}. Moreover
$(1+C_\rho)u+uD$ tends to zero. It is eventually smaller than the
fixed safe radii, the action-coincidence radius and one, while its
fine-coordinate component times the Lie-coordinate norm cost is
at most half the Haar-chart radius. Also $b<u$ and
$\tau\le\delta$ eventually.

The coefficients $m_W,t_W,W,a_*,K_5$ are bounded by fixed
polynomials in the auxiliary size parameter. The stationary jet
price has the more precise form
\[
                         p=p_0+p_1P_F,
\]
with both $p_0,p_1$ bounded by such polynomials: only the loop
price contributes the linear monitor factor. Also
$-N\log m_H$ is bounded above by a polynomial in that parameter,
including its logarithmic coupling cost.
The two action remainders
$WB_{\rm loop}$ and
$\beta W[K_5u^5+C'\tau_c^3+C'\tau_c^4]$
tend to zero, for $\tau_c\le2\tau$.
It follows that every positive term in $Y_*+\Delta+e_*+2$ is
$o(P_F^2)$. Since $\Delta=\gamma d_0^2P_F^2/4$, both inequalities
in \eqref{eq:exterior-two-reserves} hold eventually. Finally
$m_Wu\to0$ and $\Delta/\beta\to0$, proving
\eqref{eq:exterior-wire-reserve} uniformly for $\|x\|<u$.
All these thresholds can be enlarged together at a fixed stage.

We can now use \eqref{eq:exterior-exact-interior}: slack propagation
has verified its cutoff condition on the entire retained region.
Split the positive full integral into that region and its
complement, and use \eqref{eq:exterior-pointwise} and
\eqref{eq:exterior-two-reserves}. The result is
\[
 0\le Z(y)-\Omega e^{-\Phi(y)}R(b,y)
       \le\Omega e^{-\Phi(y)}\exp(-\gamma D^2/8)
\]
for every $\|y\|<u$ with $\|y-x\|\le r$, where the center $x$
is coarse healthy. This is the exterior estimate required in
Section~\ref{sec:pure-action-small-field}, on its full microscopic
comparison ball. Its hypotheses retain the actual funded sector
expansion, the stationary chart and Hessian estimates, and the
classical matching for the same coefficient data.
